\section{Pengenalan Pengujian Aplikasi Web}

Pengujian (testing) aplikasi web bertujuan memastikan fungsionalitas, kompatibilitas, dan performa sesuai harapan \cite{mdn-learn}. Jenis pengujian meliputi: \textbf{fungsional} (apakah fitur berjalan benar), \textbf{kompatibilitas} (berbagai browser dan perangkat), dan \textbf{performa} (waktu muat, responsivitas) \cite{mdn-learn}. Pengujian dapat dilakukan secara manual (menggunakan browser dan alat pengembang) atau otomatis (skrip, framework testing) \cite{webdev}.

Alat dasar: DevTools browser (tab Network, Console, Performance), validator HTML/CSS, dan untuk performa: PageSpeed Insights, Lighthouse \cite{mdn-learn}. Sub-CPMK 4.1 menuntut mahasiswa mampu melakukan pengujian dan evaluasi performa aplikasi web.
